% !TEX root = ../notes.tex

\documentclass[../notes.tex]{subfiles}

\begin{document}

\section{April 15}
Here we go.

\subsection{Some Hecke Algebras}
We continue in the setting of a reductive group $G$ over a global field $F$. Then $G(\AA_F^\infty)$ acts on the space $\mc A$ of automorphic forms. Over the finite places $v$, we have already explained that this gives rise to the action of a local Hecke algebra $\mc H_v\coloneqq C_c^\infty(G(F_v))\,dg$. Let's assemble these Hecke algebras.
\begin{definition}[restricted tensor product]
	Let $\{A_i\}$ be a collection of algebras, and select an element $e_i\in A_i$ which is idempotent for all but finitely many $i$. Then we define the \textit{restricted tensor product}
	\[\bigotimes_i(A_i,e_i)\]
	to consist of finite linear combinations of pure tensors $\bigotimes_ia_i$, where $a_i\in A_i$ for each $i$, and $a_i=e_i$ for all but finitely many $i$.
\end{definition}
One can check that this restricted tensor product forms an algebra, but we will not bother.
\begin{definition}
	Fix a reductive group $G$ over a global field $F$. For each finite place $v$, let $e_v$ be the idempotent corresponding to $1_{G(\OO_v)}$ in $\mc H_v$. We then define the \textit{nonarchimedean Hecke algebra} as
	\[\mc H_f\coloneqq\bigotimes_v(\mc H_v,e_v),\]
	where the tensor product is taken over $e_v$s.
\end{definition}
\begin{remark}
	It turns out that $\mc H_f$ is simply $C_c^\infty(G(\AA_F^\infty))\,dg$. In one direction, one sends the pure tensor $\bigotimes_vf_v\,dg_v$ to the measure $\left(\bigotimes_vf_v\right)\,dg$. (Here, the local Haar measures $dg_v$ are chosen so that the volume of $\prod_vG(\OO_v)$ is unambiguous.)
\end{remark}
\begin{remark}
	As before, we find that smooth $G(\AA_F^\infty)$-representations are equivalent to smooth $\mc H_f$-representations.
\end{remark}
We would now like to add on the archimedean places. For simplicity, we work with $\QQ$. Here is our Hecke algebra.
\begin{definition}
	Fix a reductive (real) Lie group $G$ with Lie algebra $\mf g$ and maximal compact subgroup $K\subseteq G(\RR)$. Then we define the \textit{Hecke algebra}
	\[\mc H_{\mf g,K}\coloneqq\op{Alg}(K)\otimes_{U\mf k_\CC}U\mf g_\CC.\]
	Here, $\op{Alg}(K)$ is the subalgebra of $C(K)$ consisting of all matrix coefficients of all finite-dimensional representations of $K$. It is a right module for $\mf k$, and the algebra embedding $U\mf k_\CC\subseteq U\mf g_\CC$ makes $U\mf g_\CC$ into a left module.
\end{definition}
\begin{remark}
	It turns out that $(\mf g,K)$-modules embed fully faithfully into smooth $\mc H_{g,K}$-modules. Here, smooth means that the module $M$ is the union of the modules $eM$, where $e\in\mc H_{g,K}$ is an idempotent, where $e$ the pure tensor given by an idempotent of $\op{Alg}(K)$ times the unit in $U\mf g_\CC$. Furthermore, the embedding descends to an equivalence of admissible modules.
\end{remark}
\begin{remark}
	By the Peter--Weyl theorem, the natural map
	\[\op{Alg}(K)\subseteq L^2(K)\cong\widehat{\bigoplus_{V\in\op{IrRep}(K)}}\op{End}(V)\]
	induces an isomorphism
	\[\op{Alg}(K)=\bigoplus_{V\in\op{IrRep}(K)}\op{End}(V).\]
	Indeed, the point is that $\op{Alg}(K)$ can be checked to be exactly those functions $f\in L^2(K)$ which are $K$-finite. Then one can compute what the $K$-finite vectors on the right-hand side are. Thus, each subset $\Sigma\subseteq\op{IrRep}(K)$ induces an idempotent $e_\Sigma\coloneqq\sum_{V\in\Sigma}\id_V$.
\end{remark}
\begin{example}
	Take $K=S^1$. Then the matrix coefficients of $K$ are the integer powers of the coordinate $z\in S^1$, so $\op{Alg}(K)=\CC\left[z,1/z\right]$. With the convolution algebra structure, each character $z^n$ is idempotent.
\end{example}
\begin{remark}
	Let's say something about the multiplication on $\mc H_{\mf g,K}$. The hard part is defining the multiplication of $X\cdot\varphi$, where $X\in\mf g_\CC$ and $\varphi\in\op{Alg}(K)$. If we imagine replacing $\varphi$ by a distribution $\delta_k$, then we expect to have $X\cdot\delta_k=\delta_k\cdot\op{Ad}_kX$. In general, we approximate $\varphi$ by linear combination of $\delta$s.
\end{remark}
\begin{remark}
	Here is a more conceptual definition: $\mc H_{g,K}$ is the space of distributions on $G(\RR)$ which is supported on $K$ and locally $K$-finite. Here, being supported on $K$ means that the distribution $\mu$ has $\mu(f)=0$ whenever $f|_K=0$. This definition also makes the multiplication structure clearer. It turns out that any such distribution is a linear combination of derivatives from $\mf g_\CC$ and algebraic measures, which gives the above description of $\mc H_{\mf g,K}$.
\end{remark}
We are now ready to define our global Hecke algebra.
\begin{definition}[global Hecke algebra]
	Fix a reductive group $G$ over a global field $F$. Then we define the \textit{global Hecke algebra} as
	\[\mc H\coloneqq\bigotimes_{v\text{ infinite}}\mc H_{\mf g_v,K_v}\otimes\mc H_f,\]
	where $\mf g_v=\op{Lie}G(F_v)$ and $K_v\subseteq G(F_v)$ is a maximal compact subgroup.
\end{definition}
By construction, we see that $\mc H$ admits a natural action on $\mc A$. Note that $\mc H$ remembers enough of our representation structure so that an automorphic representation is just an irreducible admissible subquotient of $\mc A$ (viewed as an $\mathcal H$-module). Indeed, the point is that we have all the correct module structures locally, which glue to the global definition.

\subsection{The Tensor Product Theorem}
We are now ready to state our theorem.
\begin{theorem}
	Let $G$ be a reductive group over a global field $F$.
	\begin{listalph}
		\item For each $v$, choose a simple $\mc H_v$-module $M_v$, and suppose that $M_v$ admits a unique spherical line $\CC m_v\in M_v^{G(\OO_v)}$ for all but finitely many $v$. Then
		\[M\coloneqq\bigotimes_v(M_v,m_v)\]
		is an irreducible $\mc H$-module.
		\item Any admissible irreducible $\mc H$-module is isomorphic to one constructed in (a).
	\end{listalph}
\end{theorem}
\begin{proof}
	For (a), let's start with admissibility. For a finite set $S$ of places, choose an idempotent $e\in\mc H$ by choosing any idempotent $e_S\in\bigotimes_{v\in S}\mc H_v$ and then selecting $e$ to be $e_v$ at each $v\notin S$. By letting $S$ expand, we see that it is enough to show that $\dim eM$ is finite for all such $e$ (because such idempotents approximate the identity). Then
	\[eM=e_SM_S\otimes\bigotimes_{v\notin S}(\CC m_v,m_v).\]
	Notably, we are only getting $m_v$s on the right because this is our chosen unique spherical vector. Note that $e_SM_S$ is finite-dimensional because $S$ is finite (and the $M_v$s for $v\in S$ are suitably smooth), so the entire space is finite-dimensional.

	We now turn to our classification of irreducible objects. Roughly speaking, the main point is to note that $M\otimes N$ is a simple module over $A\otimes B$ when $M$ and $N$ are both simple over $A$ and $B$, respectively, and all simple modules are of this form. We require the following result.
	\begin{definition}[idempotent algebra]
		An \textit{idempotent algebra} is an algebra $A$ equipped with a filtered poset of commuting idempotents $\{e_i\}$. Here, we say that $e_i\ge e_j$ if and only if $e_ie_j=e_i$.
	\end{definition}
	\begin{definition}[smooth]
		A module $M$ over an idempotent algebra $A$ is \textit{smooth} if and only if $M=\bigcup_ie_iM$. It is \textit{admissible} if and only if further $e_iM$ is finite-dimensional for each $e_i$.
	\end{definition}
	\begin{lemma}
		Fix an idempotent algebra $A$ with idempotents $\{e_i\}$. Then a smooth module $M$ is zero or simple if and only if $e_iM$ is zero or simple of finite dimension over $e_iAe_i$ for each idempotent $e_i$.
	\end{lemma}
	\begin{proof}
		Omitted.
	\end{proof}
	For example, it follows that having two smooth simple modules $M_1$ and $M_2$ over idempotent algebras $A_1$ and $A_2$ provide a smooth simple module $M\coloneqq M_1\otimes M_2$ over the idempotent algebra $A_1\otimes A_2$: indeed, for each idempotent $e_1\otimes e_2$ of $A\otimes B$, we see that
	\[(e_1\otimes e_2)M=e_1M_1\otimes e_2M_2.\]
	We now reduce to finite-dimensional considerations to find this is zero or a simple module for $e_1A_1e_1\otimes e_2A_2e_2$.\footnote{One can basically use the same arguments as one does for representations of finite groups.}

	Conversely, we claim that all simple modules $M$ of $A_1\otimes A_2$ are a tensor product. Well, we are given that $(e_1\otimes e_2)M$ is either zero or a finite-dimensional simple module of $e_1A_1e_1\otimes e_2A_2e_2$ for each idempotent $e_1\otimes e_2$. Thus, finite-dimensional considerations again allow us to decompose
	\[(e_1\otimes e_2)M=M_1(e_1)\otimes M_2(e_2)\]
	for some modules $M_1(e_1)$ and $M_2(e_2)$ unique up to isomorphism. (When $(e_1\otimes e_2)M$ vanishes, we select these modules to both vanish.) If we choose ``larger'' idempotents, meaning that $(e_1'\otimes e_2')M$ is larger, then we find that $e_1N_1(e_1')\otimes e_2N_2(e_2')\cong N_1(e_1)\otimes N_2(e_2)$. In the application, our idempotents will form an inductive system, so we receive isomorphisms $N_i(e_i)\cong e_iN_i(e_i')$ which are unique up to scaling, but then we can correct these scalars to make all the isomorphisms compatible.

	Thus, we may take the colimit of the $N_i(e_1)$s to produce a module $N_i$ of $A_i$. One can then check that the induced map
	\[N_1\otimes N_2\to M\]
	is an isomorphism by smoothness, and one can use the lemma to show that $N_1$ and $N_2$ are irreducible by passing to finite-dimensional considerations.

	We now return to $\mc H$, which is an infinite restricted tensor product. The hard part is (b), so let's focus there. Choose an irreducible admissible module $M$ of $\mc H$. For each idempotent $e$, we see $eM$ is an irreducible finite-dimensional $e\mc He$-module; we may as well take $e$ of the form $\bigotimes_ve_v'$, where $e_v'=e_v$ for all but finitely many places $v$. We collect these bad places $v$ in a finite set $S$. Now,
	\[eM=M_S\otimes\chi,\]
	where $M_S$ is some module for $e_S\mc H_Se_S$, and $\chi$ is some one-dimensional module for the commutative unital algebra $\bigotimes'_{v\notin S}e_v\mc H_ve_v$. One recovers the spherical vectors $m_v$ by looking at the embedding $\chi\into eM$, and one can use the argument above to decompose $M_S$ into a tensor product decomposition. The argument now follows the proof of the toy case with two algebras to take a colimit of the tensor product decompositions as $S$ gets large.
\end{proof}
Next class, we will say something about $L$-functions.
\begin{example}
	Fix an automorphic representation $\pi$ of $\op{GL}_{2}$ over a global field $F$. Then we can define the $L$-function of $\pi$ as
	\[L(s,\pi)\coloneqq\int_{F^\times\backslash\AA_F^\times}\varphi\left(\begin{cases}
		a \\ & 1
	\end{cases}\right)\left|a\right|^{s-1/2}\,d^\times a,\]
	where $\varphi$ is some chosen element of $\pi$. This generalizes what one expects from holomorphic modular forms. Additionally, this has the desirable property that it unfolds into an Euler product $\prod_vL_v(s,\pi_v)$, where the factors $L(s,\pi_v)$ depends only on the isomorphism class of $\pi_v$. The Euler product factorization depends on a choice of Whittaker functional.
\end{example}

\end{document}